\section{Ethics}\label{sec:ethics}

The measurement itself is low-risk. We query public APIs and storefronts from our
own machines; we do not probe networks inside censored countries, do not recruit
or expose any human participants, and do not put real users at risk of
retaliation. This is a deliberate design advantage over network-measurement
studies that can endanger in-country volunteers. We will respect platform rate
limits and terms, throttle and cache requests to avoid burdening the services, and
release only aggregate availability data plus our code.

Nonetheless, the project has ethical dimensions worth stating.

\paragraph{Selection bias in the app basket.} Our results reflect \emph{our}
choices about which apps count as ``sensitive.'' A basket assembled by a
US-based team may over-represent Western apps (NordVPN, Signal, NYT) and miss
locally important tools that are the real targets of censorship in a given
country. This bias could make some repressive environments look freer than they
are simply because we did not think to measure the right apps. We will document
our selection criteria explicitly, draw partly on external lists (e.g., the
Citizen Lab test lists and GreatFire records) to reduce our own blind spots, and
frame findings as a lower bound on censorship rather than a complete census.

\paragraph{Attribution and the fairness of claims.} Labeling an app's absence as
``censorship'' when it was actually a licensing or sanctions decision would be
unfair to the platform and could mislead readers. Our confidence-labeling
procedure (Section~\ref{sec:eval}) is partly an ethical safeguard: we commit to
distinguishing documented, inferred, and ambiguous cases rather than overstating
certainty.

\paragraph{Dual-use considerations.} Mapping which circumvention apps are
available where is primarily useful to researchers, journalists, and users seeking
access, but the same map could in principle inform a censor about gaps in its own
enforcement. We judge the transparency benefit to substantially outweigh this
risk, since the information is derived entirely from public storefronts that
censors already control, and civil-society projects already publish similar data.
